% Part: lambda-calculus
% Chapter: lambda-definability
% Section: primitive-recursive-functions

\documentclass[../../../include/open-logic-section]{subfiles}

\begin{document}

\olfileid{lam}{ldf}{prf}
\olsection{原始递归函数是$\lambda$-可定义}

回顾一下，原始递归函数是指那些可以从基本函数 $\Zero$、$\Succ$ 和 $\Proj{n}{i}$
出发，通过复合与原始递归定义出来的函数。

\begin{lem}
  \ollabel{lem:basic}
  基本的原始递归函数 $\Zero$、$\Succ$ 和投影函数~$\Proj{n}{i}$ 都是 $\lambd$-可定义的。
\end{lem}

\begin{proof}
它们分别由以下项给出 $\lambd$-定义：
\begin{align*}
  \fn{Zero} & \ident \lambd[a][\lambd[fx][x]]\\
  \fn{Succ} & \ident \lambd[a][\lambd[fx][f (a f x)]] \\
  \fn{Proj}^n_i & \ident \lambd[x_0\dots x_{n-1}][x_i]
\end{align*}
\end{proof}

\begin{lem}
  \ollabel{lem:comp} 设 $k$ 元函数 $f$ 和 $n$ 元函数 $g_0, \dots, g_{k-1}$ 分别由项
  $F$, $G_0$, \dots, $G_k$ $\lambd$-定义，而 $h$ 是由它们通过复合定义得到的函数。
  则 $h$ 是 $\lambd$-可定义的。
\end{lem}

\begin{proof}
  $h$ 可由项
  \[
  H \ident \lambd[x_0 \dots x_{n-1}][F\, (G_0 x_0 \dots
    x_{n-1}) \dots (G_{k-1} x_0 \dots x_{n-1})]
  \]
  $\lambd$-定义。此事实的验证留作练习。
\end{proof}

\begin{prob}
  请通过证明 $H\num{n_0}\dots\num{n_{n-1}} \red \num{h(n_0, \dots,
    n_{n-1})}$ 补全 \olref[lam][ldf][prf]{lem:comp} 的证明。
\end{prob}

注意，\olref{lem:comp} 并不要求 $f$ 和 $g_0$, \dots,~$g_{k-1}$ 是原始递归函数；只要求它们是全函数且 $\lambd$-可定义的。

\begin{lem}\ollabel{lem:prim}
  假设 $f$ 是一个 $n$ 元函数，$g$ 是一个 $n+2$ 元函数，
  它们分别由项 $F$ 和~$G$ $\lambd$-定义，而函数~$h$ 是由 $f$ 和 $g$ 经原始递归定义的。
  那么 $h$ 也是 $\lambd$-可定义的。
\end{lem}

\begin{proof}
  根据定义，$h$ 为：
  \begin{align*}
    h(x_1, \dots, x_n, 0) &= f(x_1, \dots, x_n)\\
    h(x_1, \dots, x_n, y+1) & = h(x_1, \dots, x_n, y, h(x_1, \dots, x_n, y)).
  \end{align*}
  直观上讲，原始递归定义将函数 $h$ 的作用迭代 $y$ 次，并将其应用于 $f(x_1,
  \dots, x_n)$。这让人联想到 Church 数码的定义，它们也是作为迭代器定义的。

  为简单起见，我们仅针对单一额外参数~$x$ 给出定义与证明。
  函数 $h$ 由以下项给出 $\lambd$-定义：
  \begin{align*}
    H \ident & \lambd[x][\lambd[y][\fn{Snd} (y
        D \tuple{\num{0}, F x})]]
    \intertext{其中 }
    D \ident & \lambd[p][\tuple{\fn{Succ} (\fn{Fst}\, p), (G x
          (\fn{Fst}\, p) (\fn{Snd}\, p))}]
  \end{align*}
  我们维护的迭代状态是一个序对，其第一分量为当前的 $y$，第二分量为对应的~$h$ 值。
  在迭代的每一步，我们构造由新的 $y$ 值与 $h$ 值组成的序对；迭代结束后，我们返回该序对的第二分量，即为最终的 $h$ 值。
  现在证明这确实是原始递归的表示。

  我们要证明，对任意 $n$ 和 $m$，$H\,\num{n}\,\num{m} \red
  \num{h(n,m)}$。
  为此，我们首先证明，如果 $D_n \ident
  \Subst{D}{\num{n}}{x}$，那么 $D_n^m \tuple{\num{0}, F\, \num{n}} \red
  \tuple{\num{m}, \num{h(n, m)}}$。
  我们对~$m$ 进行归纳。

  若 $m=0$，我们要证 $D_n^0 \tuple{\num{0}, F\, \num{n}} \red
  \tuple{\num{0}, \num{h(n, 0)}}$。
  而 $D_n^0 \tuple{\num{0}, F\,
    \num{n}}$ 即是 $\tuple{\num{0}, F\, \num{n}}$。
  由于 $F$ $\lambd$-定义了~$f$，这归约为 $\tuple{\num{0},
    \num{f(n)}}$；又因 $f(n) = h(n, 0)$，即得 $\tuple{\num{0},
    \num{h(n,0)}}$。

  现假设 $D_n^m \tuple{\num{0}, F\, \num{n}} \red
  \tuple{\num{m}, \num{h(n, m)}}$。
  我们要证 $D_n^{m+1} \tuple{\num{0}, F\, \num{n}} \red
  \tuple{\num{m+1}, \num{h(n, m+1)}}$。
  \begin{align*}
    D_n^{m+1} \tuple{\num{0}, F\, \num{n}}
    & \ident D_n(D_n^{m} \tuple{\num{0}, F\, \num{n}})\\
    & \red D_n\,\tuple{\num{m}, \num{h(n, m)}} \text{\qquad (由归纳假设)}\\
    & \ident (\lambd[p][
      \tuple{
        \fn{Succ} (\fn{Fst}\, p), (G \, \num{n} (\fn{Fst}\, p) (\fn{Snd}\, p))
    }]) \tuple{\num{m}, \num{h(n,m)}}\\
    & \redone
    \tuple{\fn{Succ} (\fn{Fst}\, \tuple{\num{m}, \num{h(n,m)}}),\\
      & \qquad (G \, \num{n} (\fn{Fst}\, \tuple{\num{m}, \num{h(n,m)}}) (\fn{Snd}\, \tuple{\num{m}, \num{h(n,m)}}))}\\
    & \red
    \tuple{\fn{Succ}\, \num{m}, (G \, \num{n} \,\num{m} \, \num{h(n,m)})}\\
    & \red \tuple{\num{m+1}, \num{g(n, m, h(n, m))}}
  \end{align*}
  由 $g(n, m, h(n, m)) = h(n, m+1)$，即得所证。

  最后，考虑
  \begin{align*}
    H\,\num n\,\num m &\ident \lambd[x][\lambd[y][\fn{Snd} (y
        (\lambd[p].\tuple{\fn{Succ} (\fn{Fst}\, p), (G\, x\,
          (\fn{Fst}\, p)\, (\fn{Snd}\, p))}) \tuple{\num{0}, F x})]]\\
    & \qquad\,\num n\, \num m\\
    & \red \fn{Snd} (\num m \,
    \underbrace{(\lambd[p].\tuple{\fn{Succ} (\fn{Fst}\, p), (G \,\num n\,
      (\fn{Fst}\, p) (\fn{Snd}\, p))})}_{D_n} \tuple{\num{0}, F \num n})\\
    & \ident \fn{Snd} (\num m \, D_n\, \tuple{\num{0}, F \num n})\\
    & \red \fn{Snd} \,(D_n^m  \tuple{\num{0}, F \num n})\\
    & \red \fn{Snd} \,\tuple{\num{m}, \num{h(n, m)}} \\
    & \red \num{h(n,m)}.
  \end{align*}
\end{proof}

\begin{prop}
  每个原始递归函数都是 $\lambd$-可定义的。
\end{prop}

\begin{proof}
  由 \olref{lem:basic}，所有基本函数都是 $\lambd$-可定义的；
  而由 \olref{lem:comp} 和 \olref{lem:prim}，$\lambd$-可定义的函数在复合与原始递归下是封闭的。
\end{proof}

\end{document}